\chapter{Finite Fourier uniformity and size bounds}
\label{core:chapter}

This chapter expands Sections~2 and~6 and Appendix~A of
\href{https://arxiv.org/abs/2110.13447}{the Ortega--Prendiville paper}.
The estimates below use explicit constants to
make the dependencies of the asymptotic notation inspectable.  The new
results in this chapter are informal proof specifications; they are not
marked as Lean-verified.  The interval estimate attributed to Cilleruelo
is proved here from a finite variance identity, so the sharper Fourier
estimate has no remaining external interval-discrepancy prerequisite.

\section{Finite sums, correlations, and the averaging kernel}

\begin{definition}[Intervals and finite Fourier sums]
\label{core:finite-fourier}
Let $N$ be a positive integer and write
$J_N=\{n\in\mathbb Z:1\leq n\leq N\}$.
All cardinalities in real inequalities are coerced to $\mathbb R$.
A finitely supported function $f:\mathbb Z\to\mathbb C$ has transform
\[
 \widehat f(\alpha)=\sum_{x\in\mathbb Z} f(x)e(\alpha x),
 \qquad e(t)=\exp(2\pi i t),\qquad \alpha\in\mathbb R/\mathbb Z.
\]
Each displayed sum over $\mathbb Z$ below has finitely many nonzero
terms.  In particular, rearrangements never require an infinite-series
convergence argument.  For two such functions put
\[
 C_{f,g}(h)=\sum_x f(x)\overline{g(x+h)},\qquad C_f=C_{f,f}.
\]
For a finite $S\subset\mathbb Z$, put
$r_S(h)=\#\{(x,y)\in S\times S:y-x=h\}$ and
$D_S=\{y-x:x,y\in S\}$.
Thus $C_{1_S}(h)=r_S(h)$, with the integer on the right regarded as a
real or complex number as appropriate.
\end{definition}

\begin{lemma}[Difference uniqueness]
\label{core:difference-uniqueness}
\uses{def:sidon_set,core:finite-fourier}
Let $S$ be a Sidon subset of an additive commutative group.  If
$x,y,u,v\in S$, $x\ne y$, and $y-x=v-u$, then $x=u$ and $y=v$.
Consequently, for a finite Sidon set,
\[
 r_S(0)=|S|,\qquad
 r_S(h)=1_{D_S}(h)\quad(h\ne0).
\]
\end{lemma}
\begin{proof}
The difference equality gives $y+u=v+x$.  The Sidon property gives
$\{y,u\}=\{v,x\}$.  Since $y$ belongs to the set on the right, either
$y=v$ or $y=x$.  The latter alternative contradicts $x\ne y$, so
$y=v$.  Cancellation in the original equality then gives $x=u$.
The pairs with difference zero are exactly $(x,x)$ with $x\in S$;
the map $x\mapsto(x,x)$ is a bijection onto that fibre.  A nonzero
fibre has at most one member by the argument just given, and has at
least one member precisely when its index lies in $D_S$.
Conversely, nonzero-difference uniqueness implies the Sidon property:
from $a+b=c+d$, either $a=c$ and cancellation gives $b=d$, or
$a-c=d-b\ne0$ and difference uniqueness gives $a=d$, $c=b$.
This also records the equivalence of the paper's definition with the
repository's definition.
\end{proof}

\begin{definition}[Integer Fej\'er averaging kernel]
\label{core:fejer-definition}
For an integer $H\geq1$, define
\[
 \mu_H(h)=\frac{\#\{(u,v)\in J_H^2:v-u=h\}}{H^2}
          =\frac{\max(H-|h|,0)}{H^2}.
\]
The first formula is the preferred definition for finite-sum proofs.
If a real parameter $T\geq1$ is used, the paper's $\mu_T$ is exactly
this kernel with $H=\lfloor T\rfloor$.
\end{definition}

\begin{lemma}[Kernel identities]
\label{core:fejer-identities}
\uses{core:fejer-definition}
For $H\geq1$, the kernel is nonnegative and even, is zero unless
$|h|<H$, and satisfies
\[
 \mu_H(0)=\frac1H,\qquad \sum_h\mu_H(h)=1,
\]
\[
 \sum_h\mu_H(h)|h|\leq H,\qquad
 \sum_h\mu_H(h)\sqrt{|h|}\leq\sqrt H.
\]
For any function $A$ on the difference interval,
\[
 \sum_h\mu_H(h)A(h)
   =\frac1{H^2}\sum_{u=1}^H\sum_{v=1}^H A(v-u).
\]
\end{lemma}
\begin{proof}
For $0\leq h<H$, the choices are $u=1,\ldots,H-h$ with
$v=u+h$, so there are $H-h$ pairs.  If $h\geq H$ there are none.
The involution $(u,v)\mapsto(v,u)$ treats negative $h$ and proves
evenness.  For $h=0$ the $H$ pairs are $(u,u)$, giving $1/H$.
Partition $J_H^2$ into the disjoint fibres of $(u,v)\mapsto v-u$;
there are $H^2$ pairs in total.  Summing over those fibres proves
both the mass-one identity and the weighted-sum identity.  On the
support, $|h|<H$, so $|h|\leq H$ and $\sqrt{|h|}\leq\sqrt H$.
Multiply these pointwise inequalities by the nonnegative weights
and use their total mass one.  These arguments also apply when
$H=1$: the kernel then has its sole nonzero value at zero.
\end{proof}

\begin{lemma}[Van der Corput, with a real-valued right side]
\lean{Sidon.vanDerCorput,Sidon.vanDerCorput_norm,Sidon.fourierSum_sq_le}\leanok
\label{core:van-der-corput}
\uses{core:finite-fourier,core:fejer-identities}
Let $N,H\geq1$ be integers and let $\operatorname{supp}(f)\subset J_N$.
Then
\[
 \left|\sum_x f(x)\right|^2
 \leq (N+H)\sum_h\mu_H(h)\operatorname{Re} C_f(h).
\]
The weighted correlation sum is real and nonnegative.  No hypothesis
$H\leq N$ is needed.  For every $\alpha\in\mathbb R/\mathbb Z$ one also
has
\[
 |\widehat f(\alpha)|^2
 \leq (N+H)\sum_h\mu_H(h)|C_f(h)|.
\]
\end{lemma}
\begin{proof}
Set $G(x)=H^{-1}\sum_{u=1}^H f(x+u)$.  For every fixed $u$, the
translation $x\mapsto x+u$ is a bijection of $\mathbb Z$, so
$\sum_xG(x)=\sum_xf(x)$.  If $G(x)\ne0$, at least one summand
is nonzero.  For its index $u$ one has $1\leq x+u\leq N$,
and therefore $1-H\leq x\leq N-1$.  This integer interval has
$N+H-1$ elements.  Complex Cauchy--Schwarz on that interval gives
\[
 |\sum_xG(x)|^2\leq (N+H-1)\sum_x|G(x)|^2
                         \leq (N+H)\sum_x|G(x)|^2.
\]
Expand the square and interchange finite sums.  In a term with
indices $u,v$, set $y=x+u$ and $h=v-u$.  This gives
\[
 \sum_x|G(x)|^2
 =\frac1{H^2}\sum_{u,v\in J_H}\sum_y
                  f(y)\overline{f(y+v-u)}
 =\sum_h\mu_H(h)C_f(h).
\]
The left side is real and nonnegative; taking real parts gives
the asserted real identity and the first inequality.  Equivalently,
$C_f(-h)=\overline{C_f(h)}$ and the kernel's evenness pair the
nonzero indices, while $C_f(0)=\sum_x|f(x)|^2$ is real.

For the Fourier inequality apply the first assertion to
$f_\alpha(x)=f(x)e(\alpha x)$.  The character is well-defined on
$\mathbb R/\mathbb Z$, satisfies $|e(\alpha h)|=1$, and gives
$C_{f_\alpha}(h)=e(-\alpha h)C_f(h)$.  Hence
$\operatorname{Re} C_{f_\alpha}(h)\leq|C_f(h)|$ at every $h$.
Multiply by $\mu_H(h)\geq0$, sum, and use $N+H>0$.
\end{proof}

\section{Appendix A: an explicit interval size bound}

\begin{lemma}[A quadratic upper bound without a quadratic formula]
\lean{Sidon.le_add_sqrt_of_sq_le}\leanok
\label{core:quadratic-bound}
If $x,b,c\geq0$ and $x^2\leq bx+c$, then $x\leq b+\sqrt c$.
\end{lemma}
\begin{proof}
Suppose instead that $x>b+\sqrt c$.  Then $x>\sqrt c$ and
$x-b>\sqrt c$.  If $c>0$, multiplying these strict inequalities
gives $x(x-b)>c$.  If $c=0$, both factors are positive and yield
the same conclusion.  This contradicts $x(x-b)\leq c$.
\end{proof}

\begin{theorem}[Erd\H{o}s--Tur\'an interval bound with a fixed constant]
\lean{IsSidonSet.card_le_sqrt_add_four_mul_rpow,IsSidonSet.card_le_five_mul_sqrt}\leanok
\label{core:sidon-size}
\uses{core:difference-uniqueness,core:van-der-corput,core:quadratic-bound}
Let $a\in\mathbb Z$, $N\geq1$, and let $S\subset\{a+1,\ldots,a+N\}$
be Sidon.  Then
\[
 |S|\leq\sqrt N+4N^{1/4}\leq5\sqrt N.
\]
For $N=0$, the interval and $S$ are empty and the same statement
holds directly.
\end{theorem}
\begin{proof}
Translation by $-a$ is injective, preserves cardinality and equality
of sums, and sends the interval to $J_N$.  It is therefore enough
to treat $S\subset J_N$.  Put $m=|S|$ and fix any integer $H\geq1$.
Apply van der Corput to $1_S$.  Separate $h=0$, where $r_S(0)=m$,
from $h\ne0$, where $0\leq r_S(h)\leq1$:
\[
 m^2\leq(N+H)\sum_h\mu_H(h)r_S(h)
 \leq(N+H)\left(\frac mH+1-\frac1H\right)
 \leq\frac{N+H}{H}m+(N+H).
\]
Both coefficients on the last line are nonnegative.  The preceding
lemma gives
\[
 m\leq\frac{N+H}{H}+\sqrt{N+H}
   \leq \sqrt N+\frac H{\sqrt N}+\frac NH+1.
\]
For the second inequality, the nonnegative number
$\sqrt N+H/\sqrt N$ has square
$N+2H+H^2/N\geq N+H$; monotonicity of the square on nonnegative
reals justifies the square-root comparison.

Take $H=\lceil N^{3/4}\rceil$.  It is a positive integer and
$N^{3/4}\leq H\leq N^{3/4}+1$.  Since $N\geq1$,
\[
 \frac H{\sqrt N}\leq N^{1/4}+N^{-1/2}\leq2N^{1/4},
 \qquad \frac NH\leq N^{1/4},\qquad 1\leq N^{1/4}.
\]
Their sum gives the stated error $4N^{1/4}$.  Finally
$N^{1/4}\leq\sqrt N$, proving the coarser bound $5\sqrt N$.
In fact $H\leq N$: the integer $N$ is an upper bound for
$N^{3/4}$, and hence also for its ceiling.  The unrestricted-window
form of van der Corput does not need this extra observation.
\end{proof}

\begin{corollary}[Local bound, including the empty interval]
\label{core:local-size}
\uses{core:sidon-size}
If $S$ is Sidon and $I$ is an interval consisting of $L$ consecutive
integers, then $|S\cap I|\leq5\sqrt L$.
\end{corollary}
\begin{proof}
If $L=0$, then $I$ and $S\cap I$ are empty.  If $L\geq1$,
the subset $S\cap I$ is Sidon and the preceding theorem applied
to its containing interval gives the claim.
\end{proof}

\section{Section 2: the four correlation errors}

\begin{lemma}[Mass captured by nonzero Sidon differences]
\label{core:captured-mass}
\uses{core:difference-uniqueness,core:van-der-corput}
If $S\subset J_N$ is Sidon, $m=|S|$, and $H\geq1$, then
\[
 W:=\sum_{h\in D_S\setminus\{0\}}\mu_H(h)
 \geq\frac{m^2}{N+H}-\frac mH.
\]
\end{lemma}
\begin{proof}
The exact identity
$\sum_h\mu_H(h)r_S(h)=m/H+W$ follows by separating zero
and using difference uniqueness at all nonzero indices.
Van der Corput applied to $1_S$ yields
$m^2\leq(N+H)(m/H+W)$; divide by the positive number $N+H$.

For completeness, the convolution proof in the paper is the same
finite counting identity in a different order.  Let
$c(t)=\#\{(s,u)\in S\times J_H:s+u=t\}$.  Expanding both sides
shows
\[
 \sum_h r_S(h)r_{J_H}(h)=\sum_t c(t)^2.
\]
Indeed, the left side counts quadruples with
$s_2-s_1=u_2-u_1$, equivalently $s_1+u_2=s_2+u_1$;
pairing $(s_1,u_2)$ and $(s_2,u_1)$ identifies these with the
pairs counted on the right.  Also $\sum_t c(t)=mH$, and $c$ is
supported on $\{2,\ldots,N+H\}$, of size $N+H-1$.
Real Cauchy--Schwarz gives
$\sum_t c(t)^2\geq(mH)^2/(N+H)$.
Dividing by $H^2$ proves the same lower bound for
$\sum_h\mu_H(h)r_S(h)$, and subtracting its zero contribution
gives the result.
\end{proof}

\begin{definition}[Normalization and error terms]
\label{core:errors}
Let $S\subset J_N$ be Sidon, set
\[
 m=|S|,\quad q=\frac m{\sqrt N},\quad
 \delta=|q-1|,\quad b=\frac{m^2}{N},\quad A=|b-1|,
\]
\[
 f_1=1_S,\qquad f_2=\frac mN1_{J_N},\qquad f=f_1-f_2.
\]
For $1\leq H\leq N$ and $i,j\in\{1,2\}$, put
\[
 C_{ij}(h)=\sum_x f_i(x)f_j(x+h),\qquad
 E_{ij}=\sum_h\mu_H(h)|b-C_{ij}(h)|.
\]
The functions $f_i$ are real-valued.  In complex Fourier formulas
they are understood through the canonical embedding into $\mathbb C$.
\end{definition}

\begin{lemma}[Interval overlap and the two mixed errors]
\label{core:interval-errors}
\uses{core:errors,core:fejer-identities,core:local-size}
One has
\[
 E_{22}\leq\frac{m^2H}{N^2},\qquad
 E_{12}=E_{21}\leq\frac{5m\sqrt H}{N}.
\]
\end{lemma}
\begin{proof}
Only $|h|<H\leq N$ contributes.  When $0\leq h<N$, the
intersection $J_N\cap(J_N-h)$ is $\{1,\ldots,N-h\}$; for
negative $h$ its cardinality is $N-|h|$ by translation, or by
the substitution $x\mapsto x+h$.  Therefore
\[
 C_{22}(h)=\frac{m^2}{N^2}(N-|h|),\qquad
 |b-C_{22}(h)|=\frac{m^2}{N^2}|h|.
\]
The first bound follows from the kernel's first-moment bound.

Define $I_h=J_N\setminus(J_N-h)$.  It is empty for $h=0$.
For $0<h<N$ it is $\{N-h+1,\ldots,N\}$; for $-N<h<0$
it is $\{1,\ldots,-h\}$.  Thus it is an interval of $|h|$
integers, including the empty-interval convention at zero.
Since $S\subset J_N$,
\[
 C_{12}(h)=\frac mN\bigl(m-|S\cap I_h|\bigr),\qquad
 |b-C_{12}(h)|=\frac mN|S\cap I_h|
                        \leq\frac{5m}{N}\sqrt{|h|}.
\]
Here the absolute value can be removed because the difference is
nonnegative.  Average with $\mu_H$ and use its square-root-moment
bound.  Finally substitution $y=x+h$ gives
$C_{21}(h)=C_{12}(-h)$.  Change $h$ to $-h$ in the finite sum
and use evenness of $\mu_H$ to obtain $E_{21}=E_{12}$.
\end{proof}

\begin{lemma}[The self-correlation error and its diagonal]
\label{core:self-error}
\uses{core:captured-mass,core:errors,core:fejer-identities}
The remaining error satisfies
\[
 E_{11}\leq 2A+\frac{m^2H}{N^2}+\frac{2m}{H}.
\]
\end{lemma}
\begin{proof}
For $h\ne0$, put $r=r_S(h)\in\{0,1\}$.  If $r=1$ then
$|b-r|=A$; if $r=0$ then $|b-r|=b\leq1+A$.
Both cases are summarized by $|b-r|\leq1-r+A$.
At zero, $C_{11}(0)=m$, and $m\leq N$ because $S\subset J_N$.
Consequently $b=m^2/N\leq m$, and
$\mu_H(0)|b-C_{11}(0)|=(m-b)/H\leq m/H$.
This contribution is essential: the estimate for nonzero $h$
cannot be applied to the diagonal.

Summing the nonzero estimate, and using total mass one, gives
\[
 E_{11}\leq\frac mH+1-W+A
 \leq 1-\frac{m^2}{N+H}+\frac{2m}{H}+A.
\]
Finally
\[
 1-\frac{m^2}{N+H}
   =1-b+\frac{bH}{N+H}
   \leq A+\frac{m^2H}{N^2},
\]
because $1-b\leq|1-b|$ and $N+H\geq N>0$.
Substitute this into the preceding inequality.
\end{proof}

\begin{lemma}[A common bound before choosing the window]
\label{core:master-bound}
\uses{core:sidon-size,core:interval-errors,core:self-error,core:van-der-corput}
For every integer $1\leq H\leq N$ and every frequency $\alpha$,
\[
 |\widehat f(\alpha)|^2
 \leq24N\delta+\frac{20N^{3/2}}{H}+200\sqrt{NH}.
\]
\end{lemma}
\begin{proof}
The size bound gives $0\leq q\leq5$, so
$A=|q^2-1|=|q-1|(q+1)\leq6\delta$.
It also gives $m^2\leq25N$.  The previous two lemmas therefore
give
\[
 E_{11}+E_{22}+E_{12}+E_{21}
 \leq12\delta+\frac{50H}{N}
       +\frac{10\sqrt N}{H}+50\sqrt{H/N}.
\]
Since $0\leq H/N\leq1$, one has $H/N\leq\sqrt{H/N}$.
Thus the right side is at most
$12\delta+10\sqrt N/H+100\sqrt{H/N}$.

Expanding the real correlation of $f=f_1-f_2$ gives
$C_f=C_{11}-C_{12}-C_{21}+C_{22}$.
Insert $b-b-b+b=0$ and use the triangle inequality to bound
$|C_f(h)|$ by the sum of the four quantities
$|C_{ij}(h)-b|$.  Average this inequality with the nonnegative
kernel, apply the Fourier form of van der Corput, and use
$N+H\leq2N$.  Multiplication by $2N$ gives the asserted constants.
\end{proof}

\begin{lemma}[Rounding a power-sized window]
\lean{Sidon.floor_rpow_window_spec}\leanok
\label{core:rounded-window}
If $N\geq1$ and $0\leq\theta\leq1$, then
$H=\lfloor N^\theta\rfloor$ satisfies
\[
 1\leq H\leq N,\qquad \tfrac12N^\theta\leq H\leq N^\theta.
\]
\end{lemma}
\begin{proof}
First $1\leq N^\theta\leq N$ by monotonicity of real powers
with base at least one.  If $1\leq x<2$, then $\lfloor x\rfloor=1$
and hence $\lfloor x\rfloor\geq x/2$.  If $x\geq2$, then
$\lfloor x\rfloor>x-1\geq x/2$.
Apply these observations to $x=N^\theta$.  The upper bound
$H\leq N$ follows because $H\leq N^\theta\leq N$.
\end{proof}

\begin{theorem}[Fourier uniformity, Theorem 1.2]
\lean{Sidon.sidon_fourier_uniformity,Sidon.sidon_fourier_uniformity_apply}\leanok
\label{core:fourier-uniformity}
\uses{core:master-bound,core:rounded-window}
For every $N\geq1$ and every Sidon $S\subset J_N$,
\[
 \left\|\widehat{1_S}-\frac{|S|}{N}\widehat{1_{J_N}}\right\|_\infty
 \leq16\sqrt N
       \left(\left|\frac{|S|}{\sqrt N}-1\right|+N^{-1/6}\right)^{1/2}.
\]
\end{theorem}
\begin{proof}
Take $H=\lfloor N^{2/3}\rfloor$.  The rounding lemma gives
$H\geq\frac12N^{2/3}$ and $H\leq N^{2/3}$, so
\[
 \frac{N^{3/2}}H\leq2N^{5/6},\qquad
 \sqrt{NH}\leq N^{5/6}.
\]
The master bound consequently gives, at every frequency,
\[
 |\widehat f(\alpha)|^2
 \leq24N\delta+240N^{5/6}
 \leq240N(\delta+N^{-1/6}).
\]
Every term on the last line is nonnegative.  Taking square roots,
using $\sqrt{240}\leq16$, proves the pointwise bound.
Linearity of the finite Fourier sum identifies $\widehat f$ with
the displayed difference.  This finite sum of continuous
characters is a continuous function on the compact circle, so
its absolute value attains a finite maximum.  Taking that maximum
preserves a bound independent of $\alpha$.  Alternatively the
pointwise bound and the least-upper-bound property give the same
supremum estimate.  Empty $S$ causes no divisions by $|S|$ and is
included; $N=1$ is included by $H=1$.
\end{proof}

\section{The Cilleruelo interval estimate, with a full finite proof}

The original article invokes an interval-discrepancy theorem from
\href{https://math.colgate.edu/~integers/a11/a11.pdf}{J.~Cilleruelo,
\emph{Gaps in dense Sidon sets}, Integers 0 (2000), A11}.
The following proof supplies that input.  Its finite-window variance
is the one used by Cilleruelo.  For the lower bound we use contracted
windows, which avoid an additional estimate for boundary intervals.
Constants are not optimized.  This proof concerns intervals of integer
points; an interval contained in $J_N$ is represented by its consecutive
integer endpoints, so no real-endpoint rounding errors are hidden.

\begin{definition}[Window counts and variance]
\label{core:window-variance-definition}
For integers $N,H\geq1$ and $S\subset J_N$, write $m=|S|$,
$T=N+H-1$, and
\[
 c_H(t)=|S\cap\{t-H+1,\ldots,t\}|,\qquad
 v=\frac{Hm}{T},\qquad
 V_H=\sum_{t=1}^{T}(c_H(t)-v)^2.
\]
In the present section put
$\lambda=\max(0,1-m/\sqrt N)$ and
$R=1+\sqrt\lambda\,N^{1/8}$.
The denominator $T$ is positive because $N,H\geq1$.
\end{definition}

\begin{lemma}[Exact window variance and its Sidon bound]
\label{core:window-variance-bound}
\uses{core:window-variance-definition,core:difference-uniqueness,core:sidon-size}
For a Sidon $S\subset J_N$ one has
\[
 V_H=Hm+2\sum_{h=1}^{H-1}(H-h)r_S(h)-\frac{H^2m^2}{T},
\]
and consequently
\[
 0\leq V_H\leq 2H^2\lambda+\frac{H^3}{N}+5H\sqrt N.
\]
\end{lemma}
\begin{proof}
For a fixed $s\in S$, the indices $t$ counted by $c_H(t)$ are
exactly $s,s+1,\ldots,s+H-1$.  All these $H$ indices lie in $J_T$,
since $1\leq s\leq N$.  Interchanging the finite sums therefore
gives $\sum_{t=1}^T c_H(t)=Hm$.
For an ordered pair $s_1,s_2\in S$, the number of indices counted
by both windows is $\max(H-|s_2-s_1|,0)$: it is the size of the
intersection of their two length-$H$ intervals.  The $m$ equal
pairs contribute $Hm$ to $\sum_t c_H(t)^2$.  For every
$1\leq h<H$, there are $r_S(h)$ pairs with difference $h$ and,
by swapping the pair, the same number with difference $-h$;
each contributes $H-h$.  This proves the second-moment formula.
Expanding the square in $V_H$, using $\sum_t c_H(t)=Hm$,
$|J_T|=T$, and $v=Hm/T$, subtracts exactly $H^2m^2/T$ and
proves the identity.

Every nonzero difference has at most one representation, so
\[
 V_H\leq Hm+H(H-1)-H^2m^2/T
      \leq Hm+H^2(1-m^2/T).
\]
The finite arithmetic sum used here is
$2\sum_{h=1}^{H-1}(H-h)=H(H-1)$; at $H=1$ both sides are zero.
Put $q=m/\sqrt N\geq0$.  If $q\geq1$, then
$N-m^2\leq0=2N\lambda$.  If $q\leq1$, then
$1-q^2=(1-q)(1+q)\leq2(1-q)=2\lambda$, which again gives
$N-m^2\leq2N\lambda$.  As $T\geq N>0$ and $H-1\geq0$,
\[
 1-\frac{m^2}{T}
 =\frac{N-m^2+H-1}{T}
 \leq\frac{2N\lambda+H-1}{T}
 \leq2\lambda+\frac HN.
\]
Finally $m\leq5\sqrt N$ by the size bound.  Substitution proves
the asserted upper bound; nonnegativity follows directly from
the definition as a sum of squares.
\end{proof}

\begin{lemma}[Variance at the three-quarter window]
\label{core:optimized-variance}
\uses{core:window-variance-bound,core:rounded-window}
For $H=\lfloor N^{3/4}\rfloor$ one has
\[
 \sqrt{V_H}\leq3N^{5/8}R,\qquad
 \frac{\sqrt{V_H}}H\leq6N^{-1/8}R,\qquad
 \frac{Hm}{N}\leq5N^{1/4}.
\]
\end{lemma}
\begin{proof}
Use $H\leq N^{3/4}$ term by term in the variance bound:
\[
 V_H\leq2N^{3/2}\lambda+6N^{5/4}
       =N^{5/4}(6+2a^2),\qquad a=\sqrt\lambda\,N^{1/8}\geq0.
\]
The inequality $6+2a^2\leq9(1+a)^2$ follows because their
difference is $3+18a+7a^2\geq0$.  Taking nonnegative square
roots gives the first claim.  Divide by
$H\geq\frac12N^{3/4}>0$ for the second claim.  The last follows
from $H\leq N^{3/4}$ and $m\leq5\sqrt N$.
\end{proof}

\begin{theorem}[Cilleruelo interval estimate, Theorem 6.1]
\label{core:cilleruelo-input}
\uses{core:optimized-variance,core:window-variance-definition}
One may take the absolute constant $C_{\mathrm C}=11$ such that,
for every $N\geq1$, every Sidon $S\subset J_N$, and every interval
$I\subset J_N$ of $L$ integers,
\[
 \left||S\cap I|-\frac{L|S|}{N}\right|
 \leq C_{\mathrm C}
       \bigl(N^{1/4}+\sqrt L\,N^{-1/8}\bigr)
       \bigl(1+\sqrt{\max(0,1-|S|/\sqrt N)}\,N^{1/8}\bigr).
\]
\end{theorem}
\begin{proof}
An empty $I$ has zero discrepancy, so suppose $I$ is nonempty.
It has the form $I=\{a+1,\ldots,b\}$ with integer
$0\leq a<b\leq N$ and $L=b-a$.  Set $k=|S\cap I|$ and take
$H=\lfloor N^{3/4}\rfloor$.  In the variance notation write
$d(t)=c_H(t)-v$ and put
$B=N^{1/4}+\sqrt L\,N^{-1/8}$.  Thus $R\geq1$ and $B>0$.

For the upper bound use the expanded interval
$U=\{a+1,\ldots,b+H-1\}\subset J_T$, of size $L+H-1$.
Every $s\in S\cap I$ contributes all its $H$ indices
$s,\ldots,s+H-1$ to $\sum_{t\in U}c_H(t)$; other contributions
are nonnegative.  Hence
\[
 Hk\leq\sum_{t\in U}c_H(t)=(L+H-1)\frac{Hm}{T}
                                      +\sum_{t\in U}d(t).
\]
Subtract $HLm/N$ and divide by $H>0$.  The deterministic term
is
\[
 m\left(\frac{L+H-1}{T}-\frac LN\right)
 =\frac{m(H-1)(N-L)}{NT}\leq\frac{mH}{N},
\]
using $0\leq L\leq N$ and $T\geq N$.
Finite Cauchy--Schwarz and $U\subset J_T$ give
\[
 \frac1H\left|\sum_{t\in U}d(t)\right|
 \leq\frac{\sqrt{|U|}\sqrt{V_H}}H
 \leq6N^{-1/8}(\sqrt L+\sqrt H)R\leq6BR.
\]
Here $\sqrt{L+H-1}\leq\sqrt L+\sqrt H$ follows by squaring
nonnegative quantities, and $\sqrt H\leq N^{3/8}$.
Combining these bounds gives
$k-Lm/N\leq5N^{1/4}+6BR\leq11BR$.

For the lower bound, if $L<H-1$, use $k\geq0$ and
$Lm/N\leq Hm/N\leq5N^{1/4}\leq11BR$.
In the remaining case $L\geq H-1$, use the contracted interval
$W=\{a+H,\ldots,b\}\subset J_T$, with cardinality $L-H+1$;
it is allowed to be empty.  For each $t\in W$ the entire
window $\{t-H+1,\ldots,t\}$ lies inside $I$.  Each member of
$S\cap I$ is counted at most $H$ times, and no member outside
$I$ is counted.  Consequently
\[
 Hk\geq\sum_{t\in W}c_H(t)
        =(L-H+1)\frac{Hm}{T}+\sum_{t\in W}d(t).
\]
After subtraction and division, the magnitude of the possible
negative deterministic term is
\[
 m\left(\frac LN-\frac{L-H+1}{T}\right)
   =\frac{m(H-1)(N+L)}{NT}\leq\frac{2mH}{N}.
\]
Cauchy--Schwarz gives
$H^{-1}|\sum_Wd(t)|\leq H^{-1}\sqrt L\sqrt{V_H}
 \leq6\sqrt L\,N^{-1/8}R$.
Thus $Lm/N-k\leq10N^{1/4}+6\sqrt L\,N^{-1/8}R\leq10BR$.
Together with the upper bound this proves absolute discrepancy
at most $11BR$, which is precisely the displayed estimate.
\end{proof}

\section{Section 6: the sharper Fourier estimate}

\begin{corollary}[Short-boundary estimate, Corollary 6.2]
\label{core:cilleruelo-short-boundary}
\uses{core:cilleruelo-input,core:sidon-size,core:errors}
Use the preceding estimate with constant $C_{\mathrm C}$, and put
$K=5+2C_{\mathrm C}$.  If $I\subset J_N$ has $L\leq N^{3/4}$
integers, then
\[
 |S\cap I|\leq K\bigl(N^{1/4}+\sqrt\delta\,N^{3/8}\bigr).
\]
\end{corollary}
\begin{proof}
The main term is bounded by
$Lm/N\leq N^{3/4}(5\sqrt N)/N=5N^{1/4}$.
The inequality $L\leq N^{3/4}$ gives
$\sqrt L\,N^{-1/8}\leq N^{1/4}$.
Also $0\leq\max(0,1-q)\leq|1-q|=\delta$, so the square-root
factor in the input is at most $\sqrt\delta$.
It follows that
\[
 |S\cap I|\leq5N^{1/4}
       +2C_{\mathrm C}N^{1/4}
       +2C_{\mathrm C}\sqrt\delta\,N^{3/8}.
\]
Both coefficients on the last two types of terms are at most $K$.
This proves the claim, including $L=0$.
\end{proof}

\begin{theorem}[Improved Fourier uniformity, Theorem 6.3]
\label{core:improved-fourier-uniformity}
\uses{core:cilleruelo-short-boundary,core:self-error,core:rounded-window,core:van-der-corput}
Let $D=290+60C_{\mathrm C}$, where $C_{\mathrm C}=11$ is available
from the proved interval estimate.  Then
\[
 \left\|\widehat{1_S}-\frac{|S|}{N}\widehat{1_{J_N}}\right\|_\infty
 \leq\sqrt D\sqrt N\bigl(\delta+N^{-1/4}\bigr)^{1/2}.
\]
\end{theorem}
\begin{proof}
Take $H=\lfloor N^{3/4}\rfloor$ and $K=5+2C_{\mathrm C}$.
Every boundary interval $I_h$ for a nonzero kernel weight has
length $|h|<H\leq N^{3/4}$.  Repeat the exact mixed-correlation
identity from Lemma~\ref{core:interval-errors}, now using the short
boundary bound:
\[
 E_{12}=E_{21}
 \leq\frac{Km}{N}(N^{1/4}+\sqrt\delta\,N^{3/8})
 \leq5K\bigl(N^{-1/4}+\sqrt\delta\,N^{-1/8}\bigr).
\]
The self-correlation and interval-overlap bounds are unchanged,
so before inserting this window size van der Corput gives
\[
 |\widehat f(\alpha)|^2
 \leq24N\delta+100H+\frac{20N^{3/2}}H
             +20KN^{3/4}+20K\sqrt\delta\,N^{7/8}.
\]
Here $H\leq N^{3/4}$ and $H\geq\frac12N^{3/4}$, so the second
and third terms sum to at most $140N^{3/4}$.  Finally
\[
 \sqrt\delta\,N^{7/8}
   =\sqrt{(N\delta)N^{3/4}}
   \leq\tfrac12(N\delta+N^{3/4}),
\]
which follows by expanding the square
$(\sqrt{N\delta}-\sqrt{N^{3/4}})^2\geq0$.
Consequently
\[
 |\widehat f(\alpha)|^2
 \leq(24+10K)N\delta+(140+30K)N^{3/4}
 \leq D N(\delta+N^{-1/4}),
\]
because $D=140+30K=290+60C_{\mathrm C}$ is at least $24+10K$.
Taking nonnegative square roots and then the supremum proceeds
exactly as in Theorem~\ref{core:fourier-uniformity}.
\end{proof}

\paragraph{Dependency discipline.}
The baseline Fourier theorem depends only on finite combinatorics,
Cauchy--Schwarz, real powers, and elementary circle characters.
The sharper theorem additionally depends on the finite
window-variance argument.  The proof of that interval estimate
uses only the elementary size bound, not either Fourier
uniformity theorem, so the dependency graph is acyclic.
